This report describes a working system, not a proposal. Every number in it was
measured on the hardware it describes, and every failure it documents was one
this project actually hit.

It is organised so that the reasoning survives a burst read. Each design
decision is stated once, followed by the consequence that forced it --- not by a
restatement of the decision in longer words. Where a choice looks arbitrary, the
measurement that produced it is given.

Two conventions are worth knowing before starting.

\textbf{Machine classes, not machine names.} The pool's concrete hostnames,
addresses, tokens and credentials live outside version control. This document
refers to \emph{the entry node}, \emph{the flagship worker}, \emph{the in-use
desktop} and so on. Nothing here is a secret, but nothing here is an inventory
either.

\textbf{Observed means observed.} Log excerpts, load figures and timings are
verbatim captures. Where something is designed but not yet exercised under real
conditions, it is labelled untested --- and Chapter~\ref{ch:failure} explains
why that distinction carries more weight in this system than in most.

The single most transferable finding is in Chapter~\ref{ch:failure}. Seven
separate defects in this deployment shared one shape: the configuration parsed
cleanly, the diagnostic tools reported correct values, and the feature did
nothing. That pattern, and the verification discipline it forces, is the part of
this work most likely to be useful to someone building something else.
